% Complete relative-growth induction; no weighted-difference or general core conclusion.
\section{Ordinary positive differences grow no faster than the main box}
Let $q_X$ be a reversed formal ordinary arm and
$K_X=[\alpha+1]\prod_{t\in U}[t]$ its main interval product.
For a product $q$ of one or two such arms, possibly multiplied
by additional interval polynomials, let $K$ be the corresponding
product of their main boxes and the same additional intervals.
Put $u_j=\max(q_j-q_{j-1},0)$, with $q_{-1}=0$.

\begin{theorem}[Relative growth of the positive differences]
For all $i<j$,
\[
                         u_j K_i\le u_i K_j.             \tag{1}
\]
The positive differences have support contained in that of $K$.
Thus their density relative to the main box is nonincreasing.
In particular $u\preceq K$, and on the strictly rising side
of $q$ every positive curvature satisfies the stronger bound
\[
              q_j-2q_{j-1}+q_{j-2}\le K_j-K_{j-1}.       \tag{2}
\]
Arbitrary internal order of the formal capacities at least
two is allowed. No claim is made for a general strict core.
\end{theorem}

\begin{proof}
We use two elementary operations on nonnegative finite sequences.
Say $a/b$ decreases when $a_jb_i\le a_ib_j$ for $i<j$,
interpreting this by cross multiplication even at zeros.
Convolution by an interval $[t]$ preserves this order.
For adjacent windows this follows by expanding the new ratio
cross product: the entering right endpoint has a no larger
ratio than each coefficient of the old window, and the outgoing
left endpoint has a no smaller ratio. Summing those pairwise
inequalities gives the assertion. The equivalent full pairwise
statement is the usual $2\times2$ minor expansion for the
interval-convolution matrix, whose minors are nonnegative.

Now let a unimodal sequence $p$ have positive and negative
difference parts $u=(\Delta p)^+$ and $v=(-\Delta p)^+$.
Their supports are ordered: every positive difference is
before every negative difference. Hence $v/u$ increases.
After convolving with $[t]$, write $U=[t]u$, $V=[t]v$.
The new positive differences are $(U-V)^+$, and their ratio
to $U$ decreases, since $V/U$ increases. Explicitly,
\[
 V_iU_j\le V_jU_i
 \ \Longrightarrow\
       (U_j-V_j)^+U_i\le(U_i-V_i)^+U_j.
\]
If $u/K$ decreases, the first window-order fact gives
$U/([t]K)$ decreasing. Multiplying these two density
comparisons proves that the positive differences of $[t]p$
have decreasing density relative to $[t]K$.
Zero support endpoints cause no issue because
$(U-V)^+\le U$; the cross-multiplied transitivity implication
is valid there as well.

For a single ordinary arm the initial interval has positive
differences supported only at zero. Appending $t$ replaces
$q$ by $[t]q+z^{a+t}$. The last added unit only changes a
terminal negative difference to zero, so it leaves the positive
differences of $[t]q$ unchanged. The preceding operation
therefore proves the induction, with new main box $[t]K$.

For two arms use induction on the combined non-one count.
A rectangular factor is the same interval-window operation.
Otherwise orient the weights as $a\ge b$. The established
ordinary-product mode theorem puts its first mode strictly
before $a$. Deleting the last part $t$ of the heavier arm gives
\[
                     q_Xq_Y=[t]q_Zq_Y+z^a q_Y.
\]
The source is strictly after the first mode. Both the whole
ordinary product and the old window are unimodal, and the
window is coefficientwise no larger than the whole product.
At the first mode they agree and attain the same global
maximum. They therefore have exactly the same positive
differences: those before that mode, where the source is absent.
Apply induction and the window operation again.
Further interval factors are handled by the same operation.

Finally $u_0=K_0=1$, so (1) implies $u\preceq K$.
At a positive rising-side curvature, put
$u_0=u_{j-1}$, $u_1=u_j$, $K_0=K_{j-1}$, $K_1=K_j$.
Here $K_0>0$, $0\le u_0\le K_0$, $u_1>u_0$ and
$u_1K_0\le u_0K_1$. These inequalities force $K_1\ge K_0$
and
\[
 (u_1-u_0)K_0\le u_0(K_1-K_0)\le K_0(K_1-K_0).
\]
Cancel $K_0$ to obtain (2).
\end{proof}

\texttt{OrdinaryRelativeGrowth.lean} checks the unbounded
finite-window cross-product transport, positive-part ratio
implication, zero-endpoint transitivity and curvature consequence.
The induction's ordinary-polynomial identifications and source
mode inputs are the written proof above, using previously
established ordinary shape theorems. No end-to-end Lean
formalization of all arm objects is asserted.

The relative-growth statement alone does not give unimodality
of $2^j u_j$: a decreasing relative density can drop quickly
and later flatten. In particular no general undilated
one-descent conclusion is inferred from this theorem.
